Road traffic injuries remain a major public health concern in Colombian cities. According to the World Health Organization's global status report, more than one million people die each year in road crashes worldwide, with a disproportionate burden in middle-income countries \cite{oms2024}. In Medellín, accident records published on the open Mede data platform provide detail on location, time, event type, and victim severity; yet this information is often consulted in a fragmented way, without tools that combine maps, indicators, and projections in a single environment \cite{mina2024,parraga2024}.\\
This degree project, entitled \textit{Traffic accident mitigation system for cities with large urban areas: Medellín case study}, presents the design and implementation of ViaData — Medellín, a web platform built to explore those historical data (approximately 2014--2021) and support accident analysis across the urban territory. The solution integrates a Django backend with PostgreSQL and the PostGIS extension, and a React frontend that publicly exposes an indicator dashboard and an analytical map with point, density, and choropleth layers. For analyst and administrator roles, it also includes a predictions module, printable reports, and a conversational assistant that queries the same system API endpoints rather than inventing figures on its own.\\
Data processing follows a reproducible ETL workflow from the Mede dataset, with centralized cleaning rules and load into a georeferenced relational database. The predictive module compares transparent statistical models---linear regression, seasonality, Poisson, moving average, $\mu \pm 3\sigma$ bands, ARIMA, and SARIMA---using hold-out validation on the last three months of the analyzed range. For monthly city-level incident counts, the adopted model was SARIMA(2,1,3)(1,1,1,12), with a mean absolute percentage error of 12.6\% on the test split. The system further provides a territorial priority index, expected load by district, trends in the proportion of fatal victims, and concentration patterns by weekday and time of day.\\
Projections are intended as orientative scenarios under stability of the filtered historical pattern; they do not claim causality or estimate individual crash risk. Overall, ViaData — Medellín delivers a geospatial instrument that links open data, interactive visualization, and documented predictive analytics to strengthen territorial readings of road risk in metropolitan settings \cite{monar2025}.

\keywords{ViaData, road accidents, Mede open data, geographic information systems, time series, urban road safety.}
